\section*{September 8, 2026: Accepted card decisions and remaining townhomes}
\textit{Agent-drafted entry. Jacob approved the seven named case resolutions; Vernon timing and final site definition remain provisional.}

The card producer now selects the approved original Assessor rows for Dearborn,
Washington, and Ingleside before summing components. Dearborn becomes one 2007
house with 5,319 square feet and 929 square feet of land. Washington becomes one
2007 building with three units, 4,684 square feet of building, and 3,101 square feet
of land; its unit count is the 2025 Assessor report, not verified original occupancy.
Ingleside becomes one 2022 three-unit building with 3,645 square feet of building
and 6,300 square feet of land. Its old building and duplicated modern card are
not extra components.

Six reviewed replacement decisions suppress the combined records for Eddy, Bell,
the three Cullerton/Prairie parcels, and 31st Street while retaining thirty
individual homes with unchanged fields. The decisions are recorded in two small
source ledgers and applied in the existing producers, with assertions about
source rows, complete assessment snapshots, replacement availability, and
competing claims. They do not change the population-wide automatic tolerances.
The rebuilt Assessor project table differs in exactly the nine approved records
from the preceding verified candidate output. Its row count is unchanged; three
extra selected card rows are removed upstream. The unchanged Make checks pass.

\textbf{Vernon.} Additional inspection explains the missing land. The later
Assessor records retain all five homes across PINs 20034000310000 (two cards),
20034000320000 (two), and 20034000330000 (one). In the same 2022 assessment their
land areas are 3,606, 3,606, and 1,803 square feet, respectively: 9,015 total,
counting each parcel once. Their building areas total 11,891 square feet. The
recorded parcel polygons are nonoverlapping and total 9,243.905 square feet,
although those saved polygons come from 2019 and 2021 rather than a common year.
The paper's five-unit record uses 11,903 square feet of building but only 3,698
square feet of land, approximately the first parcel's polygon. This appears to
assign five homes to the land of only two, overstating density. Using all three
Assessor parcels would imply about 24.2 units per acre and FAR 1.32, provisionally;
final year and common-year geography still need to be settled.

The development's own marketing page, \url{https://www.atproperties.com/developments/10310/artisan-row},
reports 2020. This agrees with the 2021 assessment that still records all five
homes together. Later individual card years also report 2021, while current
listing websites disagree among 2019, 2020, 2021, and 2023. The leading working
interpretation is 2020, but a marketing year and a permit issue date are not a
verified completion certificate. No new Vernon year or density was applied.

\textbf{Broad/Pitney.} None of the three old combined identifiers appears in the
paper's frozen input. Thirteen individual homes with PINs 17293090550000 through
17293090670000 do appear. Five Broad Street homes dated 2022 are within 500 feet;
the other eight are outside. Two additional nearby Lock Street homes also appear,
but are not included in this thirteen-home count. The combined-record confusion
does not establish that any of the individually measured homes should be dropped.

\textbf{32nd Place.} The nine individual homes at 2403--2419 West 32nd Place
have known Assessor floor areas of 1,590--1,718 square feet and construction year
2007. The uncertainty concerns the older eighteen-card record, not missing floor
areas for those nine homes. Its historical polygon contains 35 other accepted
individual homes within the date window. The inspected endpoints of the nine-home
row lie about 8--9 feet outside that polygon; this alone cannot distinguish map
imprecision from a wrong correspondence. Equal townhome areas are insufficient
to remove the eighteen-card record as their duplicate.

\textbf{Fletcher.} The strict interior-point check finds five homes for six old
cards. A sixth individual point lies only 1.55 feet outside the old polygon;
two Campbell home points lie about 23 feet outside. Thus the earlier five-home
count does not demonstrate that the sixth home is absent from the data. The
actual correspondence still needs parcel evidence because several townhome
measurements differ. No sixth observation was invented and no individually
recorded home was removed.

The preferred residential candidate and review-queue builds also completed,
and an unchanged second Make check is up to date. All nine decisions propagated.
The preferred report now has 28,070 rows versus 27,947 in its previously tracked
report; that downstream output had not yet incorporated earlier upstream changes.
Its two added fields record replacement IDs and checks. The largest newly reported
land values belong to 2025 records explicitly excluded as outside the study period;
they are not newly accepted density observations. These cumulative report changes
should not be attributed solely to the nine decisions in this entry.

The paper input's checksum remains unchanged. No density regressions were run.
Older logbook dependencies were held at saved outputs solely for compilation.
